\section{Evaluation Metrics for Graph RAG}
\label{sec:gr-metrics}

Evaluating a Graph RAG system requires assessing two coupled components---the
retriever and the generator---because an end-to-end answer can fail at either
stage~\cite{es2024ragas}. Reference-free frameworks such as RAGAS were introduced
precisely to evaluate RAG pipelines along these dimensions without relying on
human-labelled ground truth, enabling faster iteration~\cite{es2024ragas}. For
global, summarization-style questions where no single gold answer exists, GraphRAG
instead measures qualities such as the comprehensiveness and diversity of the
generated response, typically via head-to-head LLM-based
comparison~\cite{edge2024graphrag}.

\subsection{Retrieval Quality}

Retrieval quality captures whether the system surfaced the evidence actually
needed to answer the query; RAGAS operationalizes this with context-oriented
metrics that assess the relevance of the retrieved material~\cite{es2024ragas}.
This dimension is decisive because the final answer is upper-bounded by what was
retrieved---missing evidence cannot be recovered downstream~\cite{gao2023ragsurvey}.

\subsection{Answer Faithfulness and Groundedness}

Faithfulness measures whether the generated answer is actually supported by the
retrieved context rather than invented, and it is a core reference-free metric in
RAGAS~\cite{es2024ragas}. It can be made fine-grained by decomposing the answer
into atomic facts and verifying each against the source, as in atomic factual
precision scoring~\cite{min2023factscore}. Since hallucination remains a pervasive
risk in generation, faithfulness is the metric most directly aligned with the
purpose of grounding~\cite{ji2023hallucination}.
